\safechapter{Checklist final antes del examen}
% ============================================================

Debes poder hacer sin apuntes todo lo siguiente:
\begin{itemize}
  \item definir sólido rígido y sólido deformable;
  \item explicar las hipótesis de elasticidad lineal;
  \item distinguir homogeneidad e isotropía;
  \item explicar Saint-Venant y superposición;
  \item definir vector tensión;
  \item distinguir $\sigma$ y $\tau$;
  \item distinguir tensión nominal y verdadera;
  \item calcular deformación longitudinal;
  \item interpretar $\varepsilon_{ij}$;
  \item utilizar $\sigma=E\varepsilon$;
  \item reconocer apoyo móvil, articulación y empotramiento;
  \item sustituir cargas distribuidas por su resultante;
  \item aplicar el método de las secciones;
  \item dibujar correctamente $N$, $V$, $M$ y $T$ positivos;
  \item distinguir estructura isostática e hiperestática;
  \item usar $N'=-q_x$, $V'=-q_y$ y $M'=V$;
  \item deducir la forma de un diagrama sin integrar;
  \item localizar máximos de $M$ mediante $V=0$;
  \item reconocer saltos provocados por cargas puntuales y momentos;
  \item comprobar signos y unidades.
\end{itemize}

\begin{resumen}
Si puedes resolver sin ayuda los diez ejercicios resueltos, el simulacro y los seis problemas propuestos, y además explicar verbalmente las preguntas rápidas, tienes una base muy sólida del Bloque I.
\end{resumen}
